\chapter{Le second degré dans toutes ses formes}
\section*{Exercice 3.1 -- Choisir une forme}
On a $f(x)=2[(x-2)^2-4]+6=2(x-2)^2-2$, puis $f(x)=2(x-1)(x-3)$. Le carré étant positif ou nul, le minimum vaut $-2$, atteint uniquement pour $x=2$.

\section*{Exercice 3.2 -- Trois discriminants}
Les discriminants sont respectivement $0$, $25$ et $-3$. La première a la racine double $3$. La deuxième donne $(-1\pm5)/4$, donc $1$ et $-3/2$. La troisième n'a pas de racine réelle ; sa forme canonique $(x+1/2)^2+3/4$ le confirme.

\section*{Exercice 3.3 -- Inéquation factorisée}
Le produit $(x+1)(x-4)$ est positif hors des racines et négatif entre elles. Le facteur $-2$ inverse son signe. L'expression est donc positive ou nulle exactement sur $[-1,4]$. Les bornes sont incluses car l'inégalité est large.

\section*{Exercice 3.4 -- Deux racines imposées}
Le théorème de factorisation donne $P(x)=a(x+2)(x-3)$ avec $a\ne0$. La condition $P(0)=12$ donne $-6a=12$, donc $a=-2$. Le polynôme cherché est $-2(x+2)(x-3)$.

\section*{Exercice 3.5 -- Paramètre}
La forme canonique donne $(x-1)^2=1-m$. Si $m<1$, il y a deux solutions $1\pm\sqrt{1-m}$. Si $m=1$, il y a une solution double $1$. Si $m>1$, il n'y a pas de solution réelle. Cette méthode évite un calcul de discriminant supplémentaire.

\section*{Exercice 3.6 -- Optimisation géométrique}
Si un côté vaut $x$ cm, l'autre vaut $14-x$ cm avec $0<x<14$. L'aire est $x(14-x)=49-(x-7)^2\le49$. L'égalité exige $x=7$ : le rectangle est un carré de côté $7$ cm. Les unités de longueur et d'aire sont distinctes.
